\chapter{The Bell That Never Rang}

\section*{Told by}
Bell-warden Hollis, a gaunt man with soot under his fingernails and a cough that never quite stops. He tends the levee bells alone. When I asked why he had no assistant, he said, "The last one asked too many questions. The bell answered one of them. Neither of us liked the answer."

\section*{The Tale}

A bell was cast for a tower in the Mistlands. It was large, bronze, etched with ward-runes so dense that the metal itself seemed to hold its breath. The smith who forged it lost his hearing on the day of the casting --- not to the noise, but to the silence that followed. He had heard something in the cooling bronze that no living ear was meant to catch.

"Ring it once," the smith said, "and the fog retreats to the edges of the valley. Ring it twice, and the dead listen --- they do not obey, but they listen, and sometimes listening is enough. Ring it thrice, and the Direwood speaks. I cannot tell you what it says. I was not brave enough to stay for the third ring."

The warden hung the bell. He polished it. He oiled the clapper. He counted the hours without sound. The villagers grew accustomed to the silence. They forgot the bell was there. They built their lives around its promise and its refusal.

When the fog rose from the river --- thick, gray, smelling of wet ash and older things --- the villagers pleaded. "Ring the bell!"

He refused. "It is not time. The fog is testing us. It wants to know what we fear. If we ring too soon, it will learn."

When a child wandered into the mist --- a girl of seven, chasing a fox that was not a fox --- they begged at the tower door. "Ring it! The dead will help us find her."

He refused. "The dead do not help. They collect. If I ring twice, they will hear. They will come. They will find her. But they will not bring her back. They will ask what she is worth. I do not know how to answer that question."

The child was never found. The villagers never forgot. But they did not stop pleading.

When the gray wraiths came --- three of them, tall as pines, silent as the bottom of a well --- the villagers fled. They took their children and their livestock and what bread they could carry. They did not look back.

The warden stayed. He climbed the tower. His cough followed him up the stairs, step by step. He did not ring the bell. He stood beside it, arms spread, and said nothing.

The wraiths climbed the tower after him. They did not touch him. They circled the bell. They waited.

He waited longer.

At dawn, the wraiths withdrew. The fog lifted. The villagers returned.

They found the warden dead, his hands clasped around the clapper, his knuckles white, his mouth open in a silence that matched the bell's. He had kept it silent even as his breath left him. Even as the wraiths whispered offers. Even as the dead girl's voice called from the mist, asking him to ring it just once so she could find her way home.

He had not answered her. That was the part the villagers did not forgive.

They buried him under the tower. The bell still hangs. No one has rung it. The dead girl's voice still calls on certain nights, when the fog rises and the fox runs. She does not ask for the bell anymore. She asks for the warden.

\section*{The Witch's Closing}
"Some warnings should never be spoken. The bell knows. So did he. But the girl --- the girl did not know. That is the difference between the living and the dead. The dead learn too late."

\section*{What Lurks in the Shadow of the Tale}

\subsection*{The Silent Bell}

A bell that rings only when its guardian dies --- or when someone takes their place. The note it produces is not a sound. It is a cessation of sound, a hole in the world where noise goes to die. Those who hear it do not remember the tone. They remember the silence afterward.

\paragraph{Tags / Mechanics:}
\begin{itemize}
\item [WARD] [SILENCE] [SACRIFICE] [BANISH]
\item To ring the bell, a character must die while touching it. The bell tolls once --- not outward, but inward --- clearing all supernatural threats within a mile (fog, wraiths, hauntings, the dead girl's voice). The bell then cracks from throat to rim. It can be used only once. Afterward, the tower is silent forever.
\item A character may attempt to ring it without dying by succeeding on a Body + Resolve test (DV 6). On success, they suffer 3 Harm (the bell takes what it is owed) but live, and the bell cracks but can be repaired --- a task requiring a month of work and a smith who has never heard the bell's silence. On failure, they die. The bell rings anyway. The tower is silent. There is no third option.
\item The bell does not care about courage. It cares about cost. If you are not willing to pay, do not touch the clapper.
\end{itemize}

\paragraph{Adventure Hook:}
The old bell-warden lies dying in the tower. The fog is rising. The dead girl has been calling for seventy years, and the villagers have stopped hearing her --- but she has not stopped calling. The warden has never told anyone the bell's secret. On his deathbed, he whispers it to the party: "Ring it and I die for nothing. Do not ring it and she calls forever. I have kept the silence for a lifetime. Now you must choose what kind of silence you will keep."

The party must decide: let the warden die in peace and become the tower's new keepers, ring the bell and lose it forever, or find a third way --- perhaps into the mist, to answer the girl themselves.

\section*{Colophon}
Levee tower, morning. The warden coughed into his sleeve. "No," he said, when I asked if he would ring it. "Not today. Probably not tomorrow. The bell can wait. The bell has all the patience in the world. I do not. That is why I am still here."

He was still there when I left. I did not look back. The bell did not ring. The girl did not call. That day, the fog was thin enough to see through.